\documentclass{subfiles}
\begin{document}
    Let \(p > 3\) be a prime number, and let \(\K = \gf(p^{n})\).
        We call elliptic curve, a curve \(\mathcal{C} \in \mathbb{P}^{2}(\K)\) defined by an equation of the form
        \[
            x_{2}^{2}x_{0} = x_{1}^{3} + a x_{1}x_{0}^{2} + bx_{0}^{3}.
        \]
        In most cases we impose the condition \(4a + 27b \ne 0\) 
        as a way to avoid roots with multiplicity greater then one.
        In most cases we will assume \(x_{0} = 1\).
        Hence, our elliptic curves will have equations of the form \(y^{2} = x^{3} + ax + b\).
        In \Cref{fig:elliptic-curves-over-R}, we show two elliptic curves when \(\K = \rset\).
        \subfile{../TikZ/Figure 1 - Example of elliptic curves over R}

    Recall we are working with discrete fields, hence the curve will have a finite amount of point.
        In more details, our elliptic curve \(\mathcal{C}\) will represent the set of point of 
        the fieds such that its equation is satisfied. That is,
        \[
            \mathcal{C} = \set{(x, y) \in \K^{2}}[y^{2} = x^{3} + ax + b] \cup \Omega,
        \]
        where \(\Omega\) represents the point at infinity.

    \paragraph{Sum of points of an elliptic curve}
    Let \(\mathcal{C}\) be an elliptic curve, and let \(P, Q \in \mathcal{C}\).
        We want to know what \(P + Q\) is.
        Differently from what one may think, the addition of points over an elliptic curve is tricky.
        Here, we present two techniques to apply under the proper conditions.
        \begin{description}
            \item [Rope rule:] to be used when \(P \ne \pm Q\),
                consists in computing the line \(l\) passing through \(P\) and \(Q\),
                intersecting \(l\) with the equation of the curve,
                and once the system of equation is solved, 
                consider the opposite of the point found\footnotemark.

            \item [Tangent rule:] to be used when \(P = Q\),
                consists in considering the tangent line to \(P\),
                which we recall is 
                \[
                    \frac{\partial \ F}{\partial \ x}(x - x_{P}) + \frac{\partial \ F }{\partial \ y}(y - y_{P}) = 0
                \],
                intersecting it with the curve and,
                once the system of equation is solved,
                consider the opposite of the point this way found.
        \end{description}
        Below a the formulas to be used to compute the sum of two point.
        \begin{description}
            \item [\(P \ne \pm Q\):]
                \[\begin{aligned}
                    x_{P + Q} & = \left\lparen 
                        \frac{y_{P} - y_{Q}}{x_{P} - x_{Q}}
                    \right\rparen^{2} - x_{P} - x_{Q} \\
                    y_{P + Q} & = \frac{y_{P} - y_{Q}}{x_{P} - x_{Q}}(x_{P} - x_{P + Q}) - y_{P}
                \end{aligned}\]

            \item [\(P = Q\):]
                \[\begin{aligned}
                    x_{P + Q} & = x_{2P} = \left\lparen
                        \frac{3x_{P}^{2} + a}{2y_{P}}
                        \right\rparen^{2} - 2x_{p} \\ 
                    y_{P + Q} & = y_{2P} = -y_{P} + \frac{3x_{P}^{2} + a}{2y_{P}}(x_{P} - x_{2P})
                \end{aligned}\]

            \item [\(P = -Q\):] the point is \(\Omega\).
        \end{description}

    \footnotetext{The reason we consider the opposite of the point found,
        has to do with the free group of divisors of the curve;
        a topic that is far from our interests.}
\end{document}
